% Fig. 3: acquisition, delivery and analysis flowchart (single column).
% Source: 02-methodology.tex, Algoritmo de monitoreo y análisis (steps 1-7).
\begin{tikzpicture}[
  font=\scriptsize,
  every node/.append style={execute at begin node={\hyphenpenalty=10000\exhyphenpenalty=10000}},
  stp/.style={draw, rounded corners=1pt, fill=gray!12, align=center,
    text width=33mm, minimum height=7mm, inner sep=2pt},
  side/.style={draw, rounded corners=1pt, fill=white, align=center,
    text width=23mm, minimum height=7mm, inner sep=2pt},
  lane/.style={draw, dashed, rounded corners=2pt, inner sep=2mm},
  ltitle/.style={font=\scriptsize\bfseries, anchor=south west, inner sep=1pt},
  arr/.style={-{Latex[length=1.5mm]}, semithick},
  ack/.style={-{Latex[length=1.5mm]}, densely dashed, semithick},
  lbl/.style={font=\tiny, inner sep=1pt},
]
% --- Node ---
\node[stp] (s1) at (0,0) {1. Adquisición continua de pares $v$, $i$ simultáneos};
\node[stp] (s2) at (0,-1.1) {2. Bloque de 1 s: calibración, $V_{\mathrm{RMS}}$, $I_{\mathrm{RMS}}$, $P$ y calidad};
\node[stp] (s3) at (0,-2.2) {3. Resumen de 5 s: actualizar $E_k$ y punto de control};
\node[stp] (s4) at (0,-3.3) {4. Encolar en FRAM y publicar por MQTT};
\node[side] (n1) at (4.1,0) {Ventana marcada inválida};
\node[side] (n2) at (4.1,-1.1) {Umbral local confirmado: LED o zumbador y evento};
\node[side] (n4) at (4.1,-3.3) {Conservar en cola y drenar al reconectar};
\draw[arr] (s1) -- (s2);
\draw[arr] (s2) -- (s3);
\draw[arr] (s3) -- (s4);
\draw[arr] (s1) -- node[lbl, above] {CRC inválido} (n1);
\draw[arr] (s2) -- (n2);
\draw[arr] (s4) -- node[lbl, above] {sin ACK} (n4);

% --- Server ---
\node[stp] (s5) at (0,-4.85) {5. Ingesta: validar esquema e identidad, almacenar y emitir ACK};
\node[stp] (s6) at (0,-5.95) {6. Análisis: $E_h$ por hora, día y semana; hora pico $h^{*}$};
\node[stp] (s7) at (0,-7.05) {7. Alertas: abierta, reconocida, resuelta; notificación remota};
\node[side] (n5) at (4.1,-4.85) {ACK sin volver a almacenar};
\draw[arr] (s4) -- (s5);
\draw[arr] (s5) -- (s6);
\draw[arr] (s6) -- (s7);
\draw[arr] (s5) -- node[lbl, above] {duplicado} (n5);
\draw[ack] (s5.west) -- ++(-0.45,0) |- node[lbl, pos=0.25, left] {ACK} (s4.west);

\begin{scope}[on background layer]
  \node[lane, fit=(s1)(s4)(n1)(n4)] (zN) {};
  \node[lane, fit=(s5)(s7)(n5)] (zS) {};
\end{scope}
\node[ltitle] at (zN.north west) {Nodo};
\node[ltitle] at (zS.north west) {Servidor};
\end{tikzpicture}
